\section{Historia y ecosistema de la formalización matemática}

\subsection{De Grecia a la computadora: generaciones de la formalización}
La demostración matemática comenzó en Grecia; Euclides inauguró la axiomatización e Hilbert llevó la formalización al lenguaje. De la simbolización a la era del cómputo, las matemáticas se apoyaron cada vez más en la verificación automatizada: la demostración por inducción de Gentzen, la correspondencia de Curry-Howard, Coq e Isabelle se acercan a la intersección entre «legible por humanos» y «verificable por máquina».
\begin{itemize}
    \item 500 a.C.: la geometría griega traza la red axiom-theorem;
    \item 1900s: el programa de Hilbert y el sistema axiomático formal;
    \item 1970s: Curry-Howard hace que el tipo sea la demostración;
    \item 2010s: Lean y Mathlib adoptan la colaboración abierta, con la meta de que «cualquiera pueda formalize»;
    \item 2020s: AlphaProof conecta un motor de búsqueda de RL a Lean, para abordar proposiciones de nivel IMO.
\end{itemize}

Hubert también subraya que la formalización matemática no es una historia unidireccional, sino una «actualización parallel de múltiples generaciones»: cada PR de Mathlib se replay en el entorno de RL, y el RL agent a su vez retroalimenta a Mathlib con el trace de sus fallas, mientras que los governers forman discusiones interpretables en GitHub/CI; esta retroalimentación cruzada impulsa la adopción de nuevos axiomas.

\subsection{Lean, Mathlib y la estandarización}
\begin{importantbox}{Los tres principios de Lean + Mathlib}
1) \textbf{Reproducibilidad}: cada teorema cuenta con un script compilable; 2) \textbf{Unificación}: cada concepto entra en una sola biblioteca, para evitar la fragmentación de definiciones; 3) \textbf{Apertura}: Mathlib convierte a los colaboradores en nodos de gobernanza mediante PR/CI. Hubert subraya: el RL agent debe entender Mathlib y también ser sensible a los fallos de CI.
\end{importantbox}

Estos principios se reflejan directamente en el despliegue de AlphaProof: el formalizer se alinea con el pipeline de PR/CI de Mathlib, y cada vez que el agent genera un nuevo lemma debe pasar primero por el nightly runner; cuando el CI failed, el agent degrada automáticamente su search policy y, en su lugar, comparte el trace con el grupo human-in-the-loop.

Elementos semánticos que Lean ofrece a AlphaProof:
\begin{enumerate}
    \item \textbf{state del entorno}: el teorema actual, las premisas y la pila de tactic;
    \item \textbf{catalog de acciones}: los distintos tactic conforman el action space, que puede personalizarse con BFS, MCTS o Heuristic;
    \item \textbf{Retroalimentación}: la demostración exitosa equivale a +1, y una failure trigger un trace/log;
    \item \textbf{Observation}: Lean regresa el proof term, el árbol de goal y las restricciones de type.
\end{enumerate}

\subsection{Mecanismos de colaboración en el ecosistema de formalización}
Hubert describe el ciclo cerrado de colaboración de la comunidad de Lean: los mantenedores de Mathlib controlan la calidad mediante el pipeline de PR/build, y el equipo de AlphaProof sube «proof search logs + failed tactics» al governance board, lo que forma la entrada del trace de RL. Varios Lean expert abren el replay log para corregir hallucination de manera human-in-the-loop.

\begin{knowledgebox}{Puntos clave de la formalización colaborativa}
1) Crowdsourced replay log: registra la sequence de fallas del agent; 2) Fairness dashboard: verifica si la search está sesgada hacia algún tipo de política; 3) Human review loop: cuando el RL activa el drift gate, el expert puede rollback la knowledge base.
\end{knowledgebox}

\subsection{Resumen de la sección}
El desarrollo de la formalización matemática, los principios de estandarización de Lean/Mathlib y los mecanismos de colaboración community-driven sostienen en conjunto el contexto de AlphaProof; esta sección explica por qué conectar RL con Lean produce «creatividad verificable».
\newpage

